\ExamPart{Trắc nghiệm trả lời ngắn}{2}{Thí sinh trả lời từ câu 1 đến câu 4.}

\NQuestion Tìm số nghiệm thực phân biệt của phương trình $x^4-5x^2+4=0$.
\AnswerLine{$4$}
\SolutionBlock{Đặt $t=x^2$ với $t\ge0$. Khi đó phương trình trở thành
\[
t^2-5t+4=0\iff (t-1)(t-4)=0.
\]
Suy ra $t=1$ hoặc $t=4$, tức là
\[
x=\pm1,\ \pm2.
\]
Vậy phương trình có $4$ nghiệm thực phân biệt.}

\NQuestion Một xưởng in sản xuất áo đồng phục. Nếu in $x$ chiếc áo thì tổng chi phí là $C(x)=x^2-30x+400$ (với đơn vị là nghìn đồng), còn doanh thu là $R(x)=20x$
(nghìn đồng). Hỏi xưởng phải in ít nhất bao nhiêu chiếc áo để không bị lỗ?
\AnswerLine{$20$}
\SolutionBlock{Xưởng không bị lỗ khi
\[
R(x)\ge C(x).
\]
Suy ra
\[
20x\ge x^2-30x+400
\iff x^2-50x+400\le0.
\]
Ta có
\[
\Delta=50^2-4\cdot400=900\Rightarrow \sqrt{\Delta}=30.
\]
Hai nghiệm là
\[
x=\dfrac{50-30}{2}=20,\qquad x=\dfrac{50+30}{2}=40.
\]
Vì hệ số $a>0$, bất phương trình đúng khi
\[
20\le x\le40.
\]
Số áo ít nhất phải in là $20$.}

\NQuestion Một đội cứu hộ ở vị trí $A(1;2)$ cần đi đến bờ sông được mô hình bởi đường thẳng $d: 3x-4y+12=0$. Hỏi quãng đường ngắn nhất từ $A$ đến bờ sông bằng bao nhiêu?
\AnswerLine{$\dfrac{7}{5}$}
\SolutionBlock{Quãng đường ngắn nhất từ điểm $A$ đến đường thẳng $d$ chính là khoảng cách từ $A$ đến $d$: $d(A,d)=\dfrac{|3\cdot1-4\cdot2+12|}{\sqrt{3^2+(-4)^2}}=\dfrac{|3-8+12|}{5}=\dfrac{7}{5}.$
Vậy quãng đường ngắn nhất cần đi là $\dfrac{7}{5}$ (đơn vị độ dài).}

\NQuestion Tìm tung độ của trung điểm đoạn thẳng nối hai điểm $M(-2;5)$ và $N(6;-1)$.
\AnswerLine{$2$}
\SolutionBlock{Gọi $I$ là trung điểm của $MN$. Khi đó
\[
I\left(\dfrac{-2+6}{2};\dfrac{5+(-1)}{2}\right)=(2;2).
\]
Vì thế tung độ của trung điểm là $2$.}
